\chapter{El Período Especial}
\label{ch:collapse}

\section{El corte}

Lo que le ocurrió a Cuba entre 1990 y 1993 no tiene paralelo cercano entre países
que no estuvieran en guerra.

El bloque soviético había absorbido cerca del 85 por ciento del comercio cubano.
El CAME se disolvió en junio de 1991; la Unión Soviética siguió en diciembre. Se
acabaron los protocolos azucareros preferenciales, se acabó el petróleo
concesional, se acabaron los repuestos, se acabaron las líneas de crédito, y la
maquinaria de una economía entera ---construida a lo largo de dieciocho años según
estándares soviéticos, con insumos soviéticos, para un mercado soviético---
descubrió simultáneamente que ninguno de sus tres lados existía ya.

Las importaciones cubanas cayeron de unos 8.100 millones de dólares en 1989 a
unos 2.200 millones en 1993, una caída de aproximadamente tres cuartas partes en
cuatro años.\unv{} Las importaciones de petróleo cayeron de unos 13 millones de
toneladas a unos 6.\unv{} La cifra oficial cubana de la caída del producto es del
34,8 por ciento entre 1989 y 1993, probablemente el mayor descenso en cuatro años
de la Cuba del siglo XX. Es además la cifra \emph{más baja} en circulación: la
prensa cubana informó entonces del 48 por ciento, y reconstrucciones
independientes han llegado hasta el 50 por ciento \citep{mesalago2005}. La
dispersión depende de la cuestión del tipo de cambio que planteaba la nota sobre
las fuentes y no de ninguna disputa sobre lo que físicamente ocurrió, y este
libro usa la cifra oficial porque es la más conservadora disponible y el
argumento no necesita un número mayor. Para comparar: Estados Unidos perdió cerca
del 27 por ciento del producto en la Gran Depresión, a lo largo de cuatro años,
con sus socios comerciales todavía en pie.

Fidel Castro le puso nombre en 1990: el Período Especial en Tiempo de Paz, frase
tomada de los planes de contingencia de defensa civil trazados en los años ochenta
para el caso de una guerra y un bloqueo total. El nombre era exacto. Los planes se
habían escrito para un asedio, y un asedio es lo que llegó, salvo que nadie estaba
asediando.

\section{Cómo fue aquello}

La abstracción no sirve sin la textura, y la textura es lo que los cubanos que lo
vivieron siguen usando como vara para medirlo todo.

La electricidad se fue primero, porque las termoeléctricas funcionaban con
combustible importado. Los apagones de doce a dieciséis horas diarias eran
normales en 1993, programados por barrio y después ya sin programa. Sin
electricidad no hay agua, porque el agua cubana se bombea; no hay refrigeración,
en un clima donde la comida se echa a perder en horas; no hay ascensores en los
bloques de apartamentos que habían levantado las microbrigadas; y no hay luz para
leer en un país cuyo principal activo restante era que todo el mundo sabía leer.

El transporte se fue después. Los ómnibus se detuvieron por falta de combustible y
de gomas. El Estado importó más de un millón de bicicletas chinas y montó fábricas
para ensamblar más, y La Habana ---una ciudad sin ninguna tradición ciclista---
se volvió una ciudad de bicicletas, con las lesiones y el agotamiento que eso
implica con aquel calor. Después llegaron los camellos, remolques articulados
tirados por cabezas de camión, que llevaban varios cientos de pasajeros de pie
cada uno. Los bueyes volvieron a campos que llevaban veinte años mecanizados, y el
Estado los entrenó a ellos y a los hombres que los guiaban, porque no había diésel.

Después, la comida. La libreta había sido siempre un piso y la economía estatal un
techo; ahora el techo se derrumbó sobre el piso y el piso no aguantó. La ingesta
energética media diaria cayó de unas 2.900 kilocalorías en 1989 a alrededor de
1.860 en 1993, y la ingesta de proteína se redujo casi a la mitad.\unv{} Los
adultos cubanos perdieron, en promedio, entre cinco y diez kilogramos. No hubo
hambruna en el sentido técnico ---ni muertes masivas por inanición ni colapso
localizado--- pero sí hubo una reducción involuntaria, sostenida y de alcance
nacional de la ingesta de alimentos por debajo de lo que un adulto necesita.

La prueba clínica llegó en 1992. Una epidemia de neuropatía óptica y periférica se
extendió por la isla, nublando y en algunos casos destruyendo la visión. Corrió
de 1991 a 1994 y alcanzó su pico entre 1992 y 1993, y el conteo del propio
ministerio de salud cubano fue de 50.862 casos en una población de 10,8 millones
\citep{mmwr1994}. La causa era nutricional: déficit de vitaminas del grupo B,
agravado por el tabaco y el alcohol malo, en una población que llevaba dos años
comiendo demasiado poco. Se detuvo repartiendo tabletas de vitaminas a la
población entera, lo que es a la vez una acusación sobre la situación y una
demostración de lo que un sistema de salud que llega a cada cuadra puede hacer en
un mes. No hubo muertes por ella, y menos de uno de cada mil afectados quedó con
daño permanente ---que es la misma historia que las cifras de mortalidad infantil
que vienen más abajo, y hay que leerla del mismo modo---.

Todo lo demás se siguió de esas cuatro cosas. Desaparecieron el jabón y el
detergente, y con ellos llegó un aumento medible de las enfermedades diarreicas y
de la piel. La incidencia de tuberculosis se triplicó aproximadamente desde una
base muy baja. Reapareció la prostitución, dirigida a los turistas, que eran la
única fuente de divisas, y con ella una palabra, jineterismo, que no había hecho
falta en treinta años. La industria azucarera ---aquello en torno a lo cual se
había construido la economía entera--- no conseguía fertilizante, combustible ni
repuestos, y la zafra cayó de 8,1 millones de toneladas en 1989--90 a 4,3 millones
en 1993 y a 3,3 millones en 1994--95. Nunca se recuperó, en ningún año, jamás.

\section{Lo que no se derrumbó}

Aquí el libro tiene que andar con cuidado, porque esta es la evidencia más fuerte
a favor del modelo cubano y quienes la encuentran incómoda la omiten
rutinariamente.

La mortalidad infantil no subió. Bajó. En 1989 era de unos 11,1 por mil nacidos
vivos; a lo largo de los peores años del Período Especial siguió descendiendo, y
en 1994 estaba por debajo de 10.\unv{} La esperanza de vida no cayó. La cobertura
de inmunización no cayó. No cerró ninguna escuela. Las universidades siguieron
abiertas, en bicicleta y a la luz de las velas. No hubo epidemia de las firmas
clásicas de mortalidad del colapso económico ---ni un salto de muertes infantiles,
ni el desplome de la esperanza de vida masculina del tipo que Rusia registró en
exactamente los mismos años a partir de exactamente el mismo choque---.

Esto no es poca cosa y no debe explicarse a la ligera. Rusia entre 1990 y 1994
perdió una proporción comparable de producto y perdió con ella unos seis años de
esperanza de vida masculina. Cuba perdió más producto y no perdió ninguno. La
diferencia es que los sistemas de salud y educación cubanos no se financiaban con
el sector que se hundía; eran un derecho preferente sobre lo que quedara,
defendido políticamente, y atendidos por gente que siguió yendo a trabajar en un
país sin transporte. Cuando este libro sostiene en el capítulo~\ref{ch:welfare}
que el Estado de bienestar es lo único que vale la pena conservar, esta es la
evidencia: un sistema universal con alcance territorial profundo es un
amortiguador de choques de un tipo que ninguna otra herramienta al alcance de un
país pobre puede sustituir.

Dos matices, ambos reales. Los sistemas se degradaron incluso donde los
indicadores de resultado se sostuvieron ---hospitales sin anestésicos, sin
película de rayos X, sin esterilizadores en funcionamiento, y médicos que
improvisaban o se arreglaban sin ellos---. Y hubo un beneficio genuinamente
perverso, documentado después y digno de mencionarse porque muestra lo extraño del
episodio: la restricción calórica forzosa y la actividad física forzosa
produjeron una caída medible de la mortalidad por diabetes y por enfermedad
cardiovascular a lo largo de los noventa, que se revirtió cuando se recuperó el
peso después de 1995 \citep{franco2013}.\unv{} Un país adelgazó porque tenía
hambre y sus muertes por enfermedad cardíaca bajaron. Eso no es una recomendación
de política. Es un recordatorio de que este período no encaja en ninguna historia
sencilla.

\section{La Opción Cero, y lo que salió de ella}

La planificación de contingencia iba más lejos de lo que llegó a aplicarse. La
Opción Cero era el plan para la pérdida total del combustible importado: sin
generación eléctrica digna de mención, transporte de tracción animal, comida
producida y distribuida a nivel de barrio a través de cocinas comunales, el país
reorganizado en células autoabastecidas. Se preparó y no hizo falta, y que se
preparara dice algo sobre cómo evaluaba la dirigencia el año 1992.

Lo que sí se aplicó, y sobrevivió a la crisis, fue la improvisación agrícola. Sin
fertilizante, plaguicida ni combustible, la agricultura cubana se vio empujada a
métodos de bajos insumos: control biológico de plagas, compostaje, tracción
animal, y sobre todo los organopónicos, huertos de canteros elevados en solares
urbanos vacíos, que para los años dos mil abastecían una parte sustancial de las
hortalizas frescas consumidas en La Habana. Esto se hizo internacionalmente
famoso y se sobredimensiona con frecuencia. Lo que demostrablemente hizo fue poner
hortalizas en las ciudades. Lo que no hizo fue alimentar al país: la dependencia
cubana de las importaciones de alimentos subió a lo largo de todo el período, no
bajó, y el argumento a favor de la agricultura de bajos insumos hay que hacerlo
por sus méritos en el capítulo~\ref{ch:land} y no sobre una leyenda de los
noventa.

\section{Agosto de 1994}

Para el verano de 1994 la moneda se había ido. El peso, oficialmente a uno por
dólar, se cambiaba en la calle a 120 o 150. Un salario estatal no compraba casi
nada. La gente se iba en cualquier cosa que flotara, y la práctica estatal de
interceptar embarcaciones había producido una serie de incidentes violentos.

El 13 de julio de 1994 un grupo de unas setenta personas secuestró un viejo
remolcador, el \emph{13 de Marzo}, y puso rumbo a Florida. Lo persiguieron
embarcaciones operadas por el Estado que lo embistieron y dirigieron mangueras de
alta presión contra la gente en cubierta, incluidos los niños que sostenían en
alto. El remolcador se hundió a siete millas de la costa. Murieron treinta y
siete personas, diez de ellas niños.\unv{} La versión del gobierno fue que el
hundimiento fue accidental y culpa de los secuestradores. No se permitió
investigación independiente. La Comisión Interamericana de Derechos Humanos halló
responsable al Estado \citep{iachr1996}. Sea cual sea la conclusión sobre el
Período Especial en general, esto ocurrió, y les ocurrió a personas que intentaban
salir de un país que no las dejaba salir.

El 5 de agosto de 1994, tras el rumor de que se permitiría zarpar a una
embarcación, miles de personas se congregaron en el Malecón habanero, tiraron
piedras, rompieron cristales y gritaron contra el gobierno. Fue la primera
manifestación antigubernamental masiva en treinta y cinco años. Se reprimió en
cuestión de horas, en parte con policía y en parte con obreros de la construcción
movilizados con palos, y el propio Castro bajó al malecón durante los hechos. El
maleconazo duró una tarde y cambió el cálculo del gobierno de forma permanente.

A los pocos días Castro anunció que quien quisiera irse por mar podía hacerlo sin
impedimento. En las cinco semanas siguientes se echaron al mar en balsas y
neumáticos entre 30.000 y 35.000 personas.\unv{} Se ahogó un número desconocido.
Estados Unidos, incapaz de absorberlas, cambió una política de treinta años: los
balseros interceptados en el mar se llevaron a Guantánamo en vez de a Florida,
donde unos 30.000 permanecieron en campamentos durante meses. Las negociaciones
que siguieron produjeron los acuerdos migratorios de septiembre de 1994 y mayo de
1995, por los cuales Estados Unidos se comprometía a admitir un mínimo de 20.000
cubanos al año por vías legales y a devolver a los interceptados en el mar,
mientras que quienes tocaran tierra podían quedarse ---la política que se conoció
como pies secos, pies mojados y que daría forma a la emigración cubana durante los
veintidós años siguientes---.

La válvula migratoria es el mecanismo al que este libro seguirá volviendo. Dos
veces antes ---Camarioca en 1965, el Mariel en 1980--- y ahora una tercera, la
presión interna se alivió abriendo la puerta, y cada vez funcionó. Es uno de los
dos instrumentos que han mantenido estable al sistema político cubano a través de
crisis que en otro caso lo habrían roto; el otro es el piso de bienestar. El
capítulo~\ref{ch:crisis} mostrará la misma válvula abriéndose una cuarta vez, a
una escala que por fin amenazó aquello que protegía.

\begin{ledger}{el Período Especial, 1990--1994}
\emph{Logrado.} Una pérdida de más de un tercio del producto nacional absorbida
sin hambruna, sin mortalidad epidémica, sin colapso del Estado y sin la caída de
la esperanza de vida que el choque idéntico produjo en Rusia. La mortalidad
infantil siguió \emph{bajando} durante todo el período. Todas las escuelas se
mantuvieron abiertas. El sistema de salud se degradó severamente y sostuvo sus
resultados, que es la evidencia más fuerte de este libro de que un sistema
universal con alcance territorial vale lo que cuesta. La agricultura urbana,
nacida de la desesperación, se convirtió en una institución duradera y
genuinamente útil.

\emph{Costo.} Un país de lectores reducido a unas 1.860 kilocalorías diarias,
perdiendo entre cinco y diez kilos por cabeza, con cincuenta mil casos de ceguera
nutricional como recibo clínico. La industria azucarera destruida y nunca
recuperada. Una década de consumo perdida. La prostitución y una economía del
dólar restauradas como hechos de la vida cubana. Treinta y siete personas, diez de
ellas niños, ahogadas frente a La Habana en un incidente que el Estado nunca ha
investigado de forma independiente. Y el momento, una tarde de agosto en el
Malecón, en que el contrato social de 1959 dejó de ser creído por aquellos para
quienes se había escrito.
\end{ledger}
